%% kapa: exam=1 round=33 year=2022 period=1 form=A subject=civil_law subject_ko=민법 seq=12 number=12
%% source: Q-Net 감정평가사 기출문제 (cst00309) artlSeq=5209997 file=raw/1008/5209997/2022년도 제33회 감정평가사 1차시험 1교시(민법,경제학,부동산학).pdf
%% conversion: kordoc 4.24.51 -> tools/gen1.py (auto); review=proofread-sonnet; flags=none
\begin{question}{12}
사기·강박에 의한 의사표시에 관한 설명으로 옳지 않은 것은? (다툼이 있으면 판례에 따름)
\begin{choices}
  \item 상대방의 기망행위로 의사결정의 동기에 관하여 착오를 일으켜 법률행위를 한 경우, 사기를 이유로 그 의사표시를 취소할 수 있다.
  \item 상대방이 불법적인 해악의 고지 없이 각서에 서명날인할 것을 강력히 요구하는 것만으로는 강박이 되지 않는다.
  \item 부작위에 의한 기망행위로도 사기에 의한 의사표시가 성립할 수 있다.
  \item 제3자에 의한 사기행위로 계약을 체결한 경우, 표의자는 먼저 그 계약을 취소하여야 제3자에 대하여 불법행위로 인한 손해배상을 청구할 수 있다.
  \item 매수인이 매도인을 기망하여 부동산을 매수한 후 제3자에게 저당권을 설정해 준 경우, 특별한 사정이 없는 한 제3자는 매수인의 기망사실에 대하여 선의로 추정된다.
\end{choices}
\end{question}
\answer{4}
